\documentclass{gescons}

\genre {Entrevista}
\author{Virginia Ruiz}
\title{Migrações Proexológicas Internacionais: Holocarma das Nações \& Acertos Grupocármicos}
\paginaurl{https://www.youtube.com/live/r2iDY1MbxRY}

\begin{document}
    \makeentrevistatitle
    \coverart{back/Virginia}

    \begin{multicols}{2}


\begin{center}
    \includegraphics[width=8cm]{articles/entrevista/mockups/Virginia}
\end{center}


\textbf{1.       Qual foi a~motivação para a~escrita da obra? Por que a~definição deste tema para publicação de um livro?}

A motivação principal foi doar meu laboratório consciencial como voluntária internacional que migrou de Madri, Espanha, a~Foz do Iguaçu, Brasil, para participar da maxiproéxis grupal. Além disso, contribuir para a~internacionalização da Conscienciologia, trazendo neoideias e~reflexões sob a~ótica do holocarma das nações e~do Universalismo.


\textbf{2.       Quais foram as principais percepções, intra e~extrafísicas, durante a~pesquisa e~a~escrita da obra? E~posterior ao lançamento?}

Vivenciei percepções marcantes, a~exemplo do investimento paciente da equipin e~equipex, projeções lúcidas sobre a~materialização do livro e~a~identificação do público-alvo de proexistas internacionais e~migrantes intermissivistas. Senti também maior compromisso com a~interassistência tarística ao redor do mundo. Após o~lançamento, percebi expansão da lucidez pensênica e~maior abertura intraconsciencial.


\textbf{3.       Qual o~maior aprendizado com a~escrita desta obra?}

O livro me mostrou que a~migração pode ser entendida como estratégia autoevolutiva capaz de potencializar a~maxiproéxis internacional. Ao mesmo tempo, evidenciou, na prática, o~significado do Universalismo por meio da interação e~convivência com conscins de diversas culturas e~idiomas, ampliando a~compreensão da responsabilidade holocármica migratória e~demandando novos posicionamentos cosmoéticos perante grupos de retrovidas e~reciclagens intraconscienciais.

\begin{pullquote}
    ``O livro me mostrou que a~migração pode ser entendida como estratégia autoevolutiva capaz de potencializar a~maxiproéxis internacional.''
\end{pullquote}

\textbf{4.       O~que poderia dizer como incentivo para que mais pesquisadores invistam na publicação de obras conscienciológicas?}

Animo todos a~escreverem sobre sua própria teática pessoal, pois esse movimento torna a~conscin exemplarista cosmoética diante do grupo intra e~extrafísico. Em última análise, final, trata-se de assumir a~responsabilidade intermissiva que nos aproxima da compléxis e~amplia nosso papel evolutivo nestes tempos de reurbex. E~você, leitor, o~que falta para compartilhar seus achados autopesquisísticos?


\begin{pullquote}
    ``Em última análise, final, trata-se de assumir a~responsabilidade intermissiva que nos aproxima da compléxis e~amplia nosso papel evolutivo nestes tempos de reurbex.''
\end{pullquote}

    
    \end{multicols}
\end{document}
